% Options for packages loaded elsewhere
\PassOptionsToPackage{unicode}{hyperref}
\PassOptionsToPackage{hyphens}{url}
\documentclass[
]{article}
\usepackage{xcolor}
\usepackage{amsmath,amssymb}
\setcounter{secnumdepth}{-\maxdimen} % remove section numbering
\usepackage{iftex}
\ifPDFTeX
  \usepackage[T1]{fontenc}
  \usepackage[utf8]{inputenc}
  \usepackage{textcomp} % provide euro and other symbols
\else % if luatex or xetex
  \usepackage{unicode-math} % this also loads fontspec
  \defaultfontfeatures{Scale=MatchLowercase}
  \defaultfontfeatures[\rmfamily]{Ligatures=TeX,Scale=1}
\fi
\usepackage{lmodern}
\ifPDFTeX\else
  % xetex/luatex font selection
\fi
% Use upquote if available, for straight quotes in verbatim environments
\IfFileExists{upquote.sty}{\usepackage{upquote}}{}
\IfFileExists{microtype.sty}{% use microtype if available
  \usepackage[]{microtype}
  \UseMicrotypeSet[protrusion]{basicmath} % disable protrusion for tt fonts
}{}
\makeatletter
\@ifundefined{KOMAClassName}{% if non-KOMA class
  \IfFileExists{parskip.sty}{%
    \usepackage{parskip}
  }{% else
    \setlength{\parindent}{0pt}
    \setlength{\parskip}{6pt plus 2pt minus 1pt}}
}{% if KOMA class
  \KOMAoptions{parskip=half}}
\makeatother
\usepackage{longtable,booktabs,array}
\usepackage{calc} % for calculating minipage widths
% Correct order of tables after \paragraph or \subparagraph
\usepackage{etoolbox}
\makeatletter
\patchcmd\longtable{\par}{\if@noskipsec\mbox{}\fi\par}{}{}
\makeatother
% Allow footnotes in longtable head/foot
\IfFileExists{footnotehyper.sty}{\usepackage{footnotehyper}}{\usepackage{footnote}}
\makesavenoteenv{longtable}
\usepackage{graphicx}
\makeatletter
\newsavebox\pandoc@box
\newcommand*\pandocbounded[1]{% scales image to fit in text height/width
  \sbox\pandoc@box{#1}%
  \Gscale@div\@tempa{\textheight}{\dimexpr\ht\pandoc@box+\dp\pandoc@box\relax}%
  \Gscale@div\@tempb{\linewidth}{\wd\pandoc@box}%
  \ifdim\@tempb\p@<\@tempa\p@\let\@tempa\@tempb\fi% select the smaller of both
  \ifdim\@tempa\p@<\p@\scalebox{\@tempa}{\usebox\pandoc@box}%
  \else\usebox{\pandoc@box}%
  \fi%
}
% Set default figure placement to htbp
\def\fps@figure{htbp}
\makeatother
\setlength{\emergencystretch}{3em} % prevent overfull lines
\providecommand{\tightlist}{%
  \setlength{\itemsep}{0pt}\setlength{\parskip}{0pt}}
\usepackage{bookmark}
\IfFileExists{xurl.sty}{\usepackage{xurl}}{} % add URL line breaks if available
\urlstyle{same}
\hypersetup{
  hidelinks,
  pdfcreator={LaTeX via pandoc}}

\author{}
\date{}

\begin{document}

\textbf{Computational Implementation of the Surya Siddhanta Across
Fourteen Chapters}

\emph{A Project Report submitted by}

\textbf{Ashutosh Nigam}

\textbf{M25AI2006}

\emph{in partial fulfillment of the requirements for the award of the
degree of}

\textbf{M.Tech.}

\includegraphics[width=1.73056in,height=1.73056in]{./media/image1.png}

\textbf{Indian Institute of Technology Jodhpur}

\textbf{Department of Computer Science and Engineering}

\emph{November 2026}

\textbf{Declaration}

I hereby declare that the work presented in this Project Report titled
\textbf{Computational Implementation of the Surya Siddhanta Across
Fourteen Chapters} submitted to the Indian Institute of Technology
Jodhpur in partial fulfillment of the requirements for the award of the
degree of \textbf{M.Tech.}, is a Bonafide record of the research work
carried out under the supervision of \textbf{Dr. Dinesh Mohan Joshi}.
The contents of this Project Report in full or in parts, have not been
submitted to, and will not be submitted by me to, any other Institute or
University in India or abroad for the award of any degree or diploma.

\textbf{Signature}

Ashutosh Nigam

M25AI2006

\textbf{Certificate}

This is to certify that the Project Report titled \textbf{Computational
Implementation of the Surya Siddhanta Across Fourteen Chapters},
submitted by \textbf{Ashutosh Nigam (M25AI2006)} to the Indian Institute
of Technology Jodhpur for the award of the degree of \textbf{M.Tech.},
is a bonafide record of the research work done by him/her under my
supervision. To the best of my knowledge, the contents of this report,
in full or in parts, have not been submitted to any other Institute or
University for the award of any degree or diploma.

\textbf{Signature}

Dr. Dinesh Mohan Joshi

\textbf{Abstract}

This project converts the computable procedures of the Sūrya Siddhānta
into a Python library and makes them available through calculation,
reading and conversational interfaces. The completed calculation layer
contains one module for each of the text's fourteen chapters. Its scope
includes mean motions, the sine table, true planetary positions, place
and time, lunar and solar eclipses, eclipse projection, conjunctions,
fixed stars, heliacal phenomena, the lunar crescent, pātas, cosmography,
astronomical instruments and the reckoning of time.

The implementation follows Burgess's English translation and uses worked
examples from Mahavir Prasad Srivastava's Hindi \emph{Vigyan Bhashya} as
checks. Tests based on verses and examples exposed errors that routine
execution missed, especially in the interpretation of inner-planet
motion, the sign of the \emph{manda} anomaly and the Ujjain time base. A
separate DE440 comparison measures the longitude path against a modern
reference under specified conventions. It is treated as an evaluation of
this implementation, not a proof of historical practice. The report
examines the software alongside the evidence for its behavior and its
remaining limits. The completed React application includes learning
paths for the Sun, Moon, five planets and Rahu, alongside a Research
Explorer. The lessons connect a date to correction stages and source
references. The explorer lets readers inspect the saved eight-body DE440
comparison across time windows and time treatments, and download the
evidence. These additions improve access to the implemented model; no
measured improvement in student learning is claimed.

\textbf{Contents}

\begin{longtable}[]{@{}
  >{\raggedright\arraybackslash}p{(\linewidth - 2\tabcolsep) * \real{0.5000}}
  >{\raggedleft\arraybackslash}p{(\linewidth - 2\tabcolsep) * \real{0.5000}}@{}}
\toprule\noalign{}
\begin{minipage}[b]{\linewidth}\raggedright
\end{minipage} & \begin{minipage}[b]{\linewidth}\raggedleft
page
\end{minipage} \\
\midrule\noalign{}
\endhead
\bottomrule\noalign{}
\endlastfoot
Abstract & i \\
Table of Contents & ii \\
1. Introduction & \\
2. Literature Survey & \\
3. Problem Definition And Objectives & \\
4. Methodology & \\
5. Results And Findings & \\
6. Summary And Future Plan Of Work & \\
References & \\
\end{longtable}

\textbf{1. Introduction}

The \emph{Sūrya Siddhānta} gives a connected set of rules for
astronomical calculation. Its fourteen chapters begin with measures of
time and mean planetary motions, then proceed through true positions,
local geometry, eclipses, conjunctions, stars, visibility, the lunar
crescent, instruments and calendar reckoning. A verbal account can leave
the meaning of a quantity or the clock used in a calculation unresolved.
A working program requires an explicit choice for each ambiguity, and
those choices appear in its output.

The software is both an astronomical calculator and a record of
interpretation. A function that returns a longitude has to identify the
epoch and the meaning of longitude; a function that returns an eclipse
must also identify the observer. The same principle applies to
explanations generated from the corpus. They should point back to a
chapter and verse so a reader can inspect the adopted rule. The path
from source to function to test leads to the application that presents
the result.

The work has two audiences. A programmer needs a stable function and a
defined input convention. A reader of the Sanskrit astronomical
tradition needs to know whether an output follows a verse, a
translator's commentary or a software convention. The same result can
serve both audiences only if it includes the calculation's provenance.
For this reason, the project uses descriptive names for chapter modules,
records verse references in code and keeps the corpus browser close to
the calculation routes. The interfaces follow those design decisions and
remain connected to the text.

\textbf{1.1 Abbreviations and terms}

The abbreviation SS refers to the Sūrya Siddhānta. JD is Julian Day, a
continuous numerical date representation. Ahargaṇa is the elapsed day
count used by the text. Jyā is the tabulated sine. Manda and śīghra name
the two principal corrections to mean planetary position. DE440 is the
Jet Propulsion Laboratory's modern planetary and lunar ephemeris used as
a reference in a separate study. MCP is the Model Context Protocol
through which an external model client can call the project's tools.

\textbf{2. Literature Survey}

\textbf{2.1 Source texts and editorial approach}

The principal source of rules and constants is Ebenezer Burgess's 1860
English translation of the \emph{Sūrya Siddhānta}. The project also uses
worked examples in Mahavir Prasad Srivastava's Hindi \emph{Sūrya
Siddhānta: Vigyan Bhashya}. Those examples give intermediate values for
specific dates and places, which are useful for checking signs, units
and the order of operations. The Hindi material has been consulted
through OCR, so a doubtful digit must be checked against the page image.
OCR can confuse the period typeface's nine and six; a disagreement is
not automatically a software defect.

The code identifies where a procedure follows the verse and where
commentary supplies an otherwise missing step. This distinction matters
for the instruments in Chapter 13 and other passages whose descriptions
do not give a complete numerical algorithm. The project implements a
defensible reading of the computable core, not a critical edition of the
Sanskrit text and not every historical practitioner's method.

Burgess's translation is not treated as a ready-made computer program. A
verse may prescribe a relation but leave the starting value, order of
correction or instrument construction implicit. Burgess's notes
sometimes supply a method; where the project follows that
reconstruction, it records the source. The second text supplies worked
arithmetic, but does not justify removing a difference between the verse
and commentary. This separation helps explain a mismatch without
changing a historically motivated function merely to agree with one
printed table. It also shows later readers where an interpretation can
be revised.

Source attribution is especially important where a printed translation,
its notes and a later commentary offer different amounts of detail. The
project's implementation manuscript classifies disagreements with worked
examples as a software defect, an intentional deviation, a source
problem or a rule not yet implemented. This scheme prevents the
historical rule from being changed simply to make every test match a
printed number. It also prevents the opposite mistake of dismissing
every mismatch as a misprint. Each classification needs evidence: a page
image for a suspected OCR error, a verse reference for a rule, and an
intermediate trace for a calculation.

\textbf{3. Problem Definition And Objectives}

\textbf{3.1 Project objectives}

This project aims to implement those procedures in a traceable form and
test them against both the text and numerical examples. The calculation
library is central to this work. The browser applications and
conversational assistant are access layers: they let a user ask for a
position, inspect a verse or explore a diagram, while chapter functions
perform the calculations. A source passage, intermediate value and
automated check make the answer easier to question and correct.

The implementation, interfaces and evaluation described in this report
were carried out within the current project batch. The project brings
together the Sūrya Siddhānta calculation library, its applications and
its numerical evaluation. Its objectives are to implement the text's
computable procedures in reusable software, allow readers to inspect the
adopted rules and intermediate values, and evaluate the resulting
longitude calculations against a stated modern reference. The report
presents the full scope and completed deliverables together.

\textbf{3.2 Project scope and completed work}

The calculation library covers the computable core of all fourteen
chapters. The repository includes a chapter module for each chapter, a
per-planet interface, panchanga functions, REST routes, browser
interfaces, a tool-using chatbot and an MCP server. Verse-anchored
tests, worked-example checks and a DE440 longitude comparison address
distinct questions: whether a rule was transcribed consistently, whether
it reproduces a printed example and how the implementation compares with
modern positions. These checks answer different questions and should not
be combined into a single claim of accuracy.

The completed computational work follows the chapter structure. Chapters
1--3 cover mean motions, true positions, direction, place and time.
Chapters 4--6 cover eclipses and their projected geometry. Chapters
7--11 address conjunctions, stars, heliacal phenomena, the Moon and
pātas. Chapters 12--14 cover cosmography, instruments and time
reckoning. The applications let a user browse verses, run a calculation
through a web route, or call a tool through the chatbot or MCP server.
The adopted constants are explicit and are not silently combined with a
modern ephemeris.

The completed deliverables also include eight guided learning paths, a
Research Explorer and a bundled Burgess corpus containing all 500
numbered verses across fourteen chapters. The true-place module
implements the manda/śīghra correction sequence, and the chatbot's tool
registry spans the chapter modules. Tests and research manuscripts
document the computational edition and its numerical evaluation. The
report states the remaining algorithmic gaps, OCR noise and questions
requiring independent review as limitations of the delivered work.

\textbf{4. Methodology}

\textbf{4.1 Time and coordinate conventions}

Julian Day is the shared numerical date representation throughout the
library. The historical day count is anchored to the Kali epoch at the
meridian through Ujjain. The project distinguishes dates interpreted in
Ujjain local mean time from dates interpreted in Universal Time. Passing
the former to a sidereal-time routine as if it were the latter rotates
local sky geometry by about 76 degrees. Ascendants, parallax, rising and
setting and the armillary sphere all depend on this convention. An
internally consistent function can still return a wrong result when the
caller supplies the wrong clock.

The text's longitudes are sidereal. The implementation uses its own
ayaṇāṁśa to obtain a tropical value when a comparison requires one.
Location-dependent functions accept latitude and longitude, but the
reference meridian in the historical model is not Greenwich. Each
numerical example should state its date, local clock and longitude
convention.

The time convention affects tests as well as user input. A modern
interface may receive a civil clock time with a UTC offset, whereas the
chapter functions reckon from the historical meridian. The chatbot tools
convert at this boundary and label the resulting zone. This keeps each
chapter function from having to infer whether a fractional day means
local civil time, Ujjain mean time or Universal Time. The DE440
evaluation has a separate reference clock, Terrestrial Time, and
therefore computes a conversion involving the Ujjain longitude and ΔT.
Reporting these clocks explicitly prevents a time shift from being
mistaken for a flaw in planetary theory.

The historical longitude is measured from a sidereal reference. A modern
DE440 position begins as a vector in a different frame and becomes a
longitude only after a stated transformation. The project's evaluation
rotates geometric geocentric vectors to the mean ecliptic and equinox of
date, then compares circular angles. The ayaṇāṁśa used for the
historical longitude is itself a model with a back-and-forth
trepidation. A residual consequently contains both planetary-model
difference and the chosen precession model. Reporting the raw sidereal
value separately from the converted tropical value makes this
distinction visible to users of the Python API and web tools.

\textbf{4.2 Calculation architecture}

The Python package has modules ch01\_mean\_motions.py through
ch14\_reckoning\_time.py. A module corresponds to a chapter and exposes
its quantities to other modules. Later computations call earlier
positional and calendrical functions instead of copying formulas. The
dependencies can be followed directly: a mean longitude from Chapter 1
enters a true longitude from Chapter 2, which can then enter an eclipse,
conjunction, visibility calculation or panchanga element.

The planets package regroups chapter functions under Sun, Moon, Mars,
Mercury, Jupiter, Venus, Saturn, Rahu and Ketu. It is a convenience
interface, not a second planetary model. Ketu is placed opposite Rahu as
an interface convention because the text does not compute it
independently. The package metadata requires Python 3.9 or later and
NumPy as the core runtime dependency. Plotting, web services and model
access add separate dependencies at the application layer.

The FastAPI application includes calculation routers and book routers
for individual chapters. This closely follows the package structure, so
a user-visible result can be traced to a chapter function. The chatbot
tool registry also calls those functions and supplies schemas for the
model. The MCP server generates its tool list from that same registry.
This shared path limits duplicated astronomical logic, although
interface formatting and input conversion still require their own tests.
The project's architecture document and api/main.py show the current
composition; older README diagrams omit some later modules and should
not be used as the sole architecture source.

The chapter modules return structured results where a calculation has
several intermediate parts. A planetary-position result, for example,
can contain mean longitude, true longitude, \emph{manda} and
\emph{śīghra} contributions, \emph{ayaṇāṁśa}, sign and \emph{nakṣatra}.
The chatbot tool serializes selected fields to JSON, adding
human-readable context while preserving the numeric values. The API
schemas and React client provide a typed boundary for transporting these
results. This helps verification: a failing total can be traced to an
intermediate term without reimplementing the whole chapter to find the
difference.

\textbf{4.3 Chapter-wise algorithm implementation}

\textbf{4.3.1 Chapter 1 mean motions}

Chapter 1 establishes measures of time, the elapsed day count
(\emph{ahargaṇa}) and mean revolutions over a \emph{mahāyuga}. The
implementation converts dates to Julian Day, computes the elapsed count
and obtains a body's mean longitude from its revolution constant and
elapsed fraction of the cycle. It also supplies calendar quantities
reused by the panchanga and later chapters. Even a small difference in
these rates accumulates across centuries.

The inner planets require a careful reading of the revolution table.
Mercury's and Venus's listed revolutions describe their
\emph{śīghroccas}, while the mean planetary place is linked to the mean
Sun. Treating those entries as ordinary mean planetary motion makes the
later \emph{śīghra} anomaly collapse. The code now represents the
distinction explicitly, and a test based on the planets' limited
elongation from the Sun guards it. Here, a domain property catches an
error that syntax and type checks cannot.

Mean longitude is a cyclical quantity, so the elapsed count and
revolution count have to be interpreted together and reduced to a
circle. The code has separate constants for the seven classical bodies,
the lunar node and apogee. The node moves retrograde and therefore needs
a sign convention distinct from ordinary prograde motion. The project's
evaluation paper tabulates the adopted revolution counts and their
implied rates, which makes secular drift measurable. The comparisons
with modern rates evaluate the chosen constants; they do not supply
inputs to the historical calculation. Chapter 1 tests also cover
calendar and unit quantities that later chapters reuse.

The epoch itself is a mathematical input, not merely a date printed in a
preface. Chapter 1's elapsed-day count determines the phase of every
mean motion. A one-day mistake in the count has a different angular
effect for the Moon, Sun and slower planets. The code therefore includes
date conversion functions and tests for epoch and calendar behavior. The
project's later local-time correction showed that a correct day count
can still be misused when a modern clock function interprets it on
another meridian. Tests of mean motion and place-dependent calculations
therefore need to check this boundary together.

\textbf{4.3.2 Chapter 2 true positions}

Chapter 2 supplies the 24-entry \emph{jyā} table at radius 3438 and the
corrections that move a body from mean to true position. The
implementation interpolates tabulated sines and applies the \emph{manda}
and \emph{śīghra} corrections in the prescribed half and full sequence.
The \emph{manda} anomaly is measured with the sign convention of the
text. Reversing it changes the equation of centre's sign and can make a
computed solar position several degrees wrong while the calculation
still runs normally.

The routines use modern inverse trigonometric and square-root
operations, so the project implements the text's model with explicit
numerical conventions. It does not reproduce every rounding step of
historical hand arithmetic. The longitude path is particularly important
because the same positions feed later eclipse, conjunction and calendar
calculations, and because it is the part evaluated against DE440.

For the five star-planets, the \emph{śīghra} correction uses a second
moving reference. The Sun supplies that reference for Mars, Jupiter and
Saturn, while Mercury and Venus have their own \emph{śīghrocca} motion.
The project's true\_longitude\_full path applies the corrections in
sequence instead of adding two independent fixed offsets. The sequence
matters because a later correction uses a position already modified by
an earlier step. The table of \emph{jyā} values supplies the historical
trigonometric base, but interpolation and inverse operations are
numerical choices in the implementation. A comparison with a modern
longitude therefore measures the complete adopted path.

The code also computes the Moon and node from their own mean motions.
Rahu is an ascending lunar node and is represented by retrograde mean
longitude; the text supplies no corresponding periodic equation. The API
therefore should not describe Rahu's value as an observed point in the
same sense as a planet's true place. Ketu is derived at the interface as
the opposite point. These distinctions matter when evaluating the node
against a modern osculating reference: a difference can reflect the
historical mean-node model and the reference definition, as well as
numerical implementation.

\textbf{4.3.3 Chapter 3 direction place and time}

Chapter 3 turns a celestial position into local quantities for an
observer. The implemented functions cover direction, gnomon shadows,
rising and setting and the \emph{lagna}. These calculations combine
position with latitude, longitude and the local clock. Tests of these
calculations must therefore check more than a body's longitude. The same
value of JD interpreted on different meridians changes sidereal time and
hence the relation between the ecliptic and the horizon.

Correcting the Ujjain time base was one of the project's most
consequential fixes. Before it, code for later parallax and instrument
calculations could appear self-consistent while using a sky rotated
relative to the observer. A horizon, equator and ecliptic diagram
clarifies why a planetary longitude alone cannot determine a local rise
time or ascendant.

The gnomon calculation connects a measured shadow with solar altitude.
The \emph{lagna} locates the ecliptic point at the horizon, and rising
or setting computations locate the horizon crossing of an object. These
are related through the observer's local frame but are not
interchangeable outputs. The tests for this module use known geometric
relations and chapter examples. The current project guide contains a
stale Chapter 3 sample calling functions that do not exist;
ch03\_direction\_place\_time.py is the source for the current public
functions. A computational edition needs to distinguish documentation
from verified code.

The chapter's local calculations connect an abstract celestial
coordinate to an observable direction. A longitude identifies a point on
the ecliptic, but a horizon crossing also depends on obliquity, local
sidereal time and latitude. The project uses the text's obliquity
convention in its historical computations. The implementation manuscript
reports that mixing the two clocks changed local geometry so strongly
that a solar-eclipse parallax example exposed the error. This example
shows why tests of astronomical software should include the reference
frame and time base, not merely a date string and an expected angle.

\textbf{4.3.4 Chapters 4 and 5 eclipses}

Chapter 4 derives lunar eclipse circumstances from the Moon, Sun and
shadow. Chapter 5 treats solar eclipses, which need observer-dependent
corrections because the Moon's apparent position changes with location.
The implementation includes solar parallax terms \emph{lambana} and
\emph{nati}. It distinguishes a geocentric computation from one that
uses the observer's place. A result should be reported with the chosen
place and date convention, because an eclipse visible at one location
need not have the same circumstances at another.

The implementation has known limits. Full three-pass \emph{asakṛt}
iteration from a worked commentary example is absent, and some
disc-diameter helpers retain fixed mean values. These details can affect
contact estimates. The chapter implementation is therefore a useful
executable reading of the adopted procedures, while any claim about
event-timing accuracy would require a separate, dedicated evaluation.

Eclipse calculations bring several scales together: the difference
between solar and lunar longitude, the node's place, apparent diameters
and, for a solar eclipse, local parallax. A small mistake in one
quantity can shift a contact or change whether a local eclipse is
predicted. Chapter 4 and Chapter 5 keep lunar and solar eclipse paths
separate because Earth's shadow and the observer's topocentric Moon are
different geometries. The worked-example suite checks intermediate
quantities such as lambana and nati against printed values. The
limitations above restrict what agreement at one example date can
establish.

The lunar and solar cases also differ in how parallax enters. For a
lunar eclipse, the Earth's shadow and Moon are central to the geometry.
For a solar eclipse, the apparent Moon depends on the observer's
position on Earth, so a local result needs latitude and longitude. The
project's route and interface layers expose both kinds of result. Its
own evaluation paper confines its DE440 statistics to planetary
longitudes and explicitly excludes eclipse timing. A degree-scale
longitude comparison over ten-day samples cannot validate minute-scale
ingress and egress predictions, so this boundary needs to remain clear.

\textbf{4.3.5 Chapter 6 eclipse projection}

Chapter 6 converts eclipse quantities into a projected geometry with
contact points and diagram data. Separating the projection from the
event calculation helps with both testing and explanation. A numerical
calculation can be checked before it is rendered, while a figure can
make the geometry easier to understand. The project provides plotting
support as an optional part of the software; users of the core numerical
library do not need it for planetary longitudes.

To interpret an eclipse figure, a reader needs its lines and discs
identified, its date and observer specified, and a statement of whether
it shows a historical-model prediction or a modern reference. A diagram
records the geometry the program drew and must be read with the adopted
eclipse assumptions. The chapter is a distinct computational layer that
turns earlier angular quantities into an observable configuration.

The projection module is particularly useful for explaining an eclipse
result to a reader who does not know the traditional terms. Contact
points, discs and relative paths provide a geometric check on signs and
directions: a numerical magnitude alone may hide a reversed axis or
misplaced shadow. The module can return diagram data separately from
plotting, so the core calculation can be tested without a graphics
backend. A projection figure can show the chapter result with a caption
naming the calculation and its assumptions.

The projection output can also help validate the calculation. The
relative path of the eclipsing discs should agree in direction with the
signs of the computed latitude and contact offsets. If a plotted path
moves the wrong way, the defect may be in coordinate conversion or
plotting rather than the event search. Separating numerical geometry
from optional Matplotlib drawing lets tests focus on the underlying
values. In an academic presentation, a projection figure can accompany a
small table of the same contact quantities, letting the reader compare
the diagram with numbers.

\textbf{4.3.6 Chapters 7 and 8 conjunctions and stars}

Chapter 7 searches for conjunctions between planets and between a planet
and a junction star. Chapter 8 contains star and \emph{nakṣatra} data,
including functions that connect a longitude to an asterism and
calculate star rise or set. A conjunction is an event defined through
changing relative position, so the implementation must search over time
rather than calculate only a single date's value. The catalogue layer
provides named objects and their coordinates for these searches.

The coordinates of the junction stars are an interpretive issue. Chapter
8 uses an ecliptic reading of the catalogue, while Chapter 13 uses a
polar-coordinate rule to place stellar hoops on the armillary sphere.
The two readings give different declinations for some stars. The project
preserves the distinction instead of forcing every use of the catalogue
through one unexamined coordinate interpretation.

Chapter 8's catalogue also links computational values to textual names.
The \emph{nakṣatra} lookup divides the circle into asterisms and
portions, while the fixed-star functions use catalogue coordinates. The
computational-edition study derives the listed \emph{dhruvaka} values
from the shares given in the verses and records the small effect of
catalogue rounding. In Chapter 13, a different construction is needed
for stellar hoops. Running these readings side by side allows a
comparison for a particular star and date without implying that the
historical coordinate question has been settled.

The nakshatra\_catalog.json data and star functions give the catalogue a
defined software representation. A verse-based star entry can be
inspected through the corpus browser while its coordinates enter a
calculation. The computational-edition manuscript notes that the
\emph{dhruvaka} values for twenty-seven stars follow a rule based on
each star's share and the accumulated minutes of preceding asterisms. An
apparent small discrepancy for Revatī is explained by rounding in the
stored catalogue. This example shows how a tabulated historical value, a
derivation and a machine representation can agree only to a stated
precision.

\textbf{4.3.7 Chapters 9 and 10 visibility and the Moon}

Chapter 9 treats heliacal rising and setting: the return of a planet or
star to visibility after it has been close to the Sun. This is not the
same event as a geometric conjunction. The library therefore provides
dedicated functions for finding the next heliacal event. Chapter 10
addresses the Moon's elongation, illuminated fraction, phase name,
crescent width, horn angle, moonrise and moonset. The calculations
combine angular position with local geometry where necessary.

These chapters show how a shared positional core supports questions that
require additional assumptions. A phase can be inferred from Sun--Moon
elongation, while a rise or set time needs a place. A claim of naked-eye
visibility is more demanding than a geometric calculation and should be
phrased according to the particular rule implemented. The chapter
modules make these distinctions clear to callers.

The Moon functions produce several kinds of output from the same
Sun--Moon geometry. Elongation is an angle, illuminated fraction is a
geometric consequence of phase, and a crescent's width and tilt depend
on further relations. Moonrise and moonset also require the observer.
The project's guide gives an end-to-end example for these Chapter 10
calls, while tests exercise the implementation. For heliacal events, the
code searches forward from a starting JD rather than requiring the user
to guess the event date. This design is suitable for an exploratory
interface, provided that the output is explained as the result of the
adopted visibility criterion.

The applications make these outputs available in several forms. A
positions page can show the Sun and Moon on a zodiac wheel; a Moon page
can show phase and crescent geometry; a star page can use a sky map.
These visuals explain the same chapter results. They also expose where
an input is missing. A phase calculation can run for a date alone, while
moonrise needs a site. The chatbot's context includes a user place and
UTC offset when available, and its tools are instructed to state what
location and clock they used.

\textbf{4.3.8 Chapter 11 patas}

Chapter 11 implements Vyatipāta and Vaidhṛta as \emph{pāta} events
associated with equal declinations. The calculation corrects lunar
declination for latitude, searches for the event and estimates its
beginning and end from the relevant motion. This chapter shows why a
name alone should not define an algorithm. The \emph{pāta} here differs
from the daily panchanga \emph{nitya-yoga} terms that share the names
Vyatipāta and Vaidhṛta.

A verse-level check changed an earlier duration computation: the
denominator is the difference of daily motions required by the passage,
not simply the rate at which two already-computed declinations separate.
This difference has a material effect on the duration. The current
module, API routes and user-interface page make the event inspectable
while tests preserve the adopted interpretation.

The Chapter 11 module is exposed through find\_patas and
compute\_pata\_details chatbot tools, calculation routes and a Patas
page. These interfaces allow checks that all outputs share the same
meaning. The event has a true instant and an interval; the interval
depends on the rate specified in the verse, so a mistaken denominator
changes duration even if the central instant looks plausible. The
chapter's separate name and tests also help prevent confusion with
panchanga \emph{yoga}, where a similarly named daily category is
computed from the sum of solar and lunar longitudes.

The \emph{pāta} calculation shows how a derived event can magnify a
small definitional mistake. The centre of the event is found by a
condition on declinations, but the start and end depend on how quickly
the relevant quantities change. A rate based on the wrong motion makes
an interval too short even if the central condition is still reached.
The implementation changelog records the verse-based correction and the
corresponding API and UI exposure. Its tests can separately check the
declination relation at the true instant and the interval arithmetic.
Checking these separately makes the numerical result easier to audit.

\textbf{4.3.9 Chapter 12 cosmography}

Chapter 12 treats the Earth's globe, distances in the text's cosmic
model and cycles of time. The software derives planetary orbit distances
from the Moon's stated 324,000 yojanas and the assumption that each body
covers the same linear distance in a yuga. It also implements the
Earth's circumference, a latitude-adjusted circumference and a longitude
difference correction needed elsewhere in the book. These are
computations within the \emph{Sūrya Siddhānta}'s model. They should not
be described as modern measured distances.

The module also represents descriptive material through a verse-tagged
catalogue. Separating numerical derivations from descriptive entries
lets an interface present the source without implying that every
cosmographic statement is a contemporary astronomical measurement. The
chapter shows how the project handles material that is partly procedural
and partly explanatory.

The cosmography calculation reproduces stated orbit values from a common
linear-distance assumption and the Chapter 1 revolution counts. The
stated Moon orbit, shared rule and derived orbit values can be compared
directly with units. The project changelog notes that the derived values
reproduce the text's figures closely and connect the Sun orbit to a
constant already used in the lunar-eclipse module. That connection is an
internal consistency check on the edition. It does not turn the yojana
values into empirical modern distances.

Chapter 12 also includes material that a web page can display without
turning it into a numeric calculation: creation order, cosmic time
divisions and other descriptive entries. The data catalogue tags these
entries by verse. The numerical functions cover Earth dimensions, orbit
derivation and time regents. In the project changelog, the derived orbit
values agree with the stated values to within about 1.3 yojanas, a
precise internal check for the implemented rule. The figure is a result
of arithmetic and rounding in the historical model; it should not be
confused with uncertainty in modern physical distances.

\textbf{4.3.10 Chapter 13 instruments}

Chapter 13 describes instruments rather than only a sequence of
planetary formulas. The implementation represents the \emph{gola-yantra}
as a parametric armillary sphere in an equatorial frame. Hoops for the
equator, ecliptic, horizon, colures, solar paths and stars can be passed
to a renderer. The React interface contains an interactive
three-dimensional representation, while the underlying module stores the
astronomical geometry independently of the graphics library.

The chapter also implements the clepsydra and time finding from the
gnomon, wheel, arc and staff. Some of these procedures are supplied by
Burgess's commentary rather than fully by the verses. The output
therefore identifies its source category. This allows a user to
distinguish an operation directly warranted by the text from a useful
reconstruction that completes an underspecified description.

The armillary model stores a hoop by its plane and offset, so a renderer
can construct a sphere without recomputing astronomy. It changes with
place and moment: horizon orientation and the ecliptic point on it are
local, while some reference circles remain fixed in the model frame. The
interactive page makes this movement visible. The instruments section
also shows why source labelling matters. Where a verse names a device
but does not give every step needed to calculate time from it, the
software may use Burgess's commentary, but the result says so. Stating
the added rule presents that uncertainty more faithfully than hiding it.

The Chapter 13 model uses an equatorial coordinate frame whose north
pole and equinox axes are fixed by convention. A hoop is a geometric
object in that frame, and a page renderer can rotate the view without
changing the underlying astronomical values. For Ujjain, the
computational-edition manuscript reports 46 hoops in its model. It also
states that the eastern-horizon ecliptic point is computed directly
instead of interpolated between zodiac signs. This matters because
interpolation could make the interactive sphere look smooth while giving
a measurably wrong local ascendant. The implementation separates the
geometry from its display.

\textbf{4.3.11 Chapter 14 reckoning of time}

Chapter 14 returns to calendrical questions. The implementation covers
measures of time, solar ingress, months and seasons, and named Jovian
years. These operations draw on earlier solar positions and the
historical day count, so they inherit the date and meridian conventions
already discussed. Their outputs often use familiar calendar labels, but
each label comes from a calculation with a particular rule and epoch.

The chapter also includes alternative ways of naming a Jovian year. The
verse rule and the mean-longitude method supplied in commentary remain
distinct where both are exposed. This follows the project's approach of
preserving competing procedures and their provenance when a single
answer would hide an interpretive decision.

The code supplies time-unit conversion through a catalogue of
\emph{mānas} and calculates solar ingress events for a Gregorian year.
It can also return ayanas, equinox-related crossings and
\emph{ṣaḍaśītimukha} divisions. The Fathers' day calculation and
Jupiter's named-year cycle have separate methods. The barhaspatya\_year
function lets a caller select a method; the result therefore needs to
identify the method as well as the year name. Chapter 14 handles
calendrical categories whose definitions depend on the textual source,
rather than treating a familiar modern calendar label as
self-explanatory.

The time-reckoning chapter's solar\_year function is organized around
the Sun's entry into Aries in a Gregorian year. Related functions return
the ingress events and seasonal divisions, while convert handles
numerical conversion between measures of time. The \emph{pitrya\_clock}
function addresses the Fathers' day using the Sun--Moon relation. These
are distinct routes through the chapter; they should not be summarized
as a single generic ``calendar converter.'' The named Jovian year also
needs a method label, because a mean-motion method in Burgess's
discussion and a method tied to heliacal behavior need not return the
same label at every date.

\textbf{4.4 Public interfaces}

Programmers can use the chapter modules and Python API directly. The
Gradio explorer groups results into pages for positions, panchanga,
eclipses, conjunctions, stars and later chapters. The separate daily
panchanga application presents a narrower almanac view. Streamlit pages
cover several calculation groups, and a React frontend calls FastAPI
routes for both calculations and book browsing. The FastAPI entry point
includes chapter-specific routers, corpus, glossary and chat routes.

All these interfaces depend on the chapter layer. Reimplementing a
formula independently for each interface would make disagreements
difficult to diagnose. A screenshot or sample response establishes what
a user can inspect when its inputs are identified. Correctness comes
from tests and source comparison, rather than the appearance of the
interface.

The React frontend includes dedicated pages for the later chapters,
including Cosmography and Instruments, which are not mirrored in every
older interface. Its verse-browser pages call book routes, while
interactive SVG and three-dimensional components render astronomical
geometry. The Streamlit and Gradio applications have their own page
sets. The applications share the same core package, but their available
capabilities differ and should be described accurately. The REST layer
makes it convenient to inspect JSON outputs and keeps the user interface
separate from chapter arithmetic. An architecture figure should place
the chapter package at the bottom, API and tool layers above it, and
applications at the top.

The applications support different research tasks. A book page lets a
reader inspect a verse and nearby text. A calculation page gives a
result for chosen inputs. A visual component can show a configuration
that a table alone obscures. The API route can preserve the precise JSON
for a test or another program. The frontend also provides glossary
access, so users do not need to know every Sanskrit term before
exploring a page. The interfaces therefore provide several ways into the
same computation, but the source code and tests remain the authoritative
record of what was calculated. Learning paths (/learn) reuse the Chapter
1 and 2 APIs; /research presents the saved evaluation data.

\textbf{4.5 Corpus retrieval and chatbot}

The chatbot can call calculation tools across all fourteen chapters and
retrieval tools that return verse-indexed text. The agent prompt directs
it to call a calculation tool for a number and to retrieve the relevant
passage before quoting or paraphrasing the book. The retrieval layer
returns identifiers and source metadata that a generated explanation can
cite. When the indexed SQLite corpus is unavailable or empty, the
retrieval code can build a local FTS5 index from a JSON export.

The repository includes a bundled Burgess corpus containing all 500
numbered verses across fourteen chapters. Chat retrieval, corpus search
and Verse Explorer use this export by default, while explicit JSON and
database overrides remain supported. Repaired verse boundaries and
source locators improve traceability, but untouched OCR can still
contain transcription noise. Optional Vigyan Bhashya commentary remains
a separate source. The hosted assistant requires credentials for its
configured provider: OpenRouter, direct DeepSeek or a custom compatible
endpoint. The MCP server exposes the same calculation and corpus tool
registry over standard input and output, with the client providing its
own model.

The retrieval functions include single-verse lookup, chapter ranges,
search and chapter index. Search uses identifiers rather than returning
anonymous text snippets, which gives the chatbot a stable way to cite a
passage. The agent also distinguishes a Burgess verse from a result
taken from Srivastava's commentary; it should not attribute a commentary
explanation to the ancient text. The tool prompt instructs the model to
take computed numbers from calculation tools, not from its memory or an
OCR passage. This is an architectural safeguard, though not a guarantee
of a correct final sentence. A complete demonstration records a
question, the calculation tool called, the retrieved verse ID and the
response.

The assistant prompt specifies a bounded tool loop and asks that
numerical answers come from the registered calculation functions. It can
retrieve Burgess verse records and, when available, Vigyan Bhashya
commentary hits, labelling the latter separately. The retrieval layer
uses an indexed SQLite database when one is supplied and populated;
otherwise, it builds a local FTS5 index from the bundled JSON export.
The checked-in Burgess corpus can therefore support citations without a
runtime dependency on the external OCR pipeline. The explanations the
assistant can support still depend on commentary availability and the
quality of the remaining OCR.

\textbf{4.6 Test strategy}

Verification has three layers. Verse-anchored tests ask whether each
function follows the adopted rule. Worked-example checks ask whether a
chain of calculations reproduces an example at its printed precision.
The DE440 comparison asks how far a particular longitude path departs
from a modern reference under specified conventions. Each layer catches
a different kind of problem. A code test can pass even when the chosen
reading of a verse is wrong; a modern comparison can be close by chance;
a printed example may contain a typographical or OCR error.

The computational-edition manuscript reports 1,256 chapter and
per-planet interface tests and 257 selected worked-example cases. It
also states that the audit-case counts refer to an audit branch whose
cases and exact revision have not yet been archived. These are reported
project results, not a newly executed count from this checkout. A future
release can record a fixed code revision and the exact test command
alongside the audit cases.

The strongest tests combine numerical assertions with source-based
properties. An elongation bound catches a semantically wrong
inner-planet formula. A printed eclipse example tests intermediate
parallax, which can catch a clock shift. Wrapper tests check that the
per-planet entry points remain aligned with their chapter functions. The
repository also includes API and chatbot tests, so the interface path
can be checked independently of the chapter mathematics. A test count
belongs to a particular date and revision because this repository is
still being developed. The description of how each class of test relates
to its sources remains valid even when the count changes.

The test organization follows the chapter structure, with dedicated
files for Chapters 1--14 and additional files for API, chatbot, corpus
and per-planet interfaces. The computational-edition manuscript
separates chapter tests from wrapper tests, and a recent DE440
revalidation ran a focused Chapter 1--2 suite of 263 tests. Those
figures refer to particular checks and revisions, so they should not be
added together as a new total. A unit test should encode a textual or
geometric expectation, while a worked-example test should preserve the
printed inputs and output precision.

\textbf{4.7 Learning and research interfaces}

\textbf{4.7.1 Mars Learning Lab}

The /learn/mars route turns the Chapter 2 true-place calculation into a
guided experiment for Mars. A whole-day input and a slider spanning 90
days on either side of the chosen date control the existing book API.
Results are evaluated at midnight in Ujjain mean time and displayed as
sidereal angular directions from Earth. The drawing shows angles, not
physical orbital distances.

\includegraphics[width=6.15in,height=3.83261in,alt={Figure 1 Mars Learning Lab showing the final correction stage for 4 November 2026}]{./media/current-report-1.png}

Figure 1 Mars Learning Lab showing the final correction stage for 4
November 2026

The five selectable positions are the mean place, half sighra, half
manda, whole manda on the original mean, and whole sighra. The fourth
step explicitly returns to the original mean, so the sequence should not
be read as a cumulative sum. At the demonstration date the final
longitude is 114.3004 degrees. A link opens the source calculation with
the same date; optional corpus text is distinguished from calculation
availability.

An understanding check asks whether the final position is ahead of or
behind the mean along the shortest angular separation. JSON export
preserves the inputs, time basis, frame and intermediate results. This
teaching functionality is implemented; its effect on student learning
remains unmeasured.

\textbf{4.7.2 Learning paths}

The /learn directory contains eight working lessons: Sun, Moon, Mars,
Mercury, Venus, Jupiter, Saturn and Rahu. Learners can start with the
Sun, then study the Moon before tracing the planetary corrections. The
directory supports direct links to any lesson; it does not require a
prescribed order or claim measured mastery.

\begin{longtable}[]{@{}
  >{\raggedright\arraybackslash}p{(\linewidth - 2\tabcolsep) * \real{0.5000}}
  >{\raggedright\arraybackslash}p{(\linewidth - 2\tabcolsep) * \real{0.5000}}@{}}
\toprule\noalign{}
\begin{minipage}[b]{\linewidth}\raggedright
Learning path
\end{minipage} & \begin{minipage}[b]{\linewidth}\raggedright
Implemented experiment
\end{minipage} \\
\midrule\noalign{}
\endhead
\bottomrule\noalign{}
\endlastfoot
Sun & Mean and manda-corrected longitude; tropical Sun, declination and
day/night length at Ujjain. \\
Moon & Two correction positions; Sun--Moon elongation, 12-degree tithi
intervals and waxing/waning paksha. \\
Mars & The original five-position lesson and its saved research
panel. \\
Mercury and Venus & Individual lessons and a same-date comparison of
every correction stage and Sun separation. \\
Jupiter and Saturn & Individual lessons and a same-date comparison of
the five-position sequence. \\
Rahu & Ascending and descending lunar nodes, regression, longitude
geometry and an eclipse concept quiz. \\
\end{longtable}

All seven body lessons share date validation, a direction diagram,
selectable API-derived stages, chapter references, a prediction check
and JSON export. Dedicated panels relate the calculation to each topic.
Sun and Moon have two positions rather than the five used for the
star-planets. Planet comparisons use one API response so both bodies
share the input date.

The seasonal panel explicitly changes from sidereal to tropical
longitude and states that it uses the library\textquotesingle s
hour-angle day-length method. The Moon panel displays elongation and
tithi without claiming to render the illuminated disc. Each tithi spans
12 degrees of longitude separation and need not last one civil day.
Paired-planet Sun separation is a longitude difference, not an
observational visibility estimate.

Downloads include the lesson\textquotesingle s additional context:
day-length values for the Sun, tithi results for the Moon, or comparison
rows for paired planets. Late responses cannot overwrite a newer date.
Invalid or failed requests keep exports disabled and label any previous
result. The numerical lessons remain available if the optional corpus
fails.

\textbf{4.7.3 Lunar nodes and eclipse geometry}

The Rahu lesson at /learn/rahu draws the ascending lunar node returned
by Chapter 1 and derives the opposite descending direction by adding 180
degrees modulo 360. The Sun and Moon positions come from Chapter 2 for
the same input date. Learners can select either node or the Moon
direction and observe regression through the node\textquotesingle s
negative mean motion.

\includegraphics[width=6.15in,height=3.075in,alt={Figure 2 Rahu lesson showing the two lunar nodes and same-date Sun and Moon directions}]{./media/current-report-2.png}

Figure 2 Rahu lesson showing the two lunar nodes and same-date Sun and
Moon directions

The projection shows ecliptic longitude only, with orbital tilt
described in text. The lesson explains that solar eclipses require new
Moon near a node and lunar eclipses require full Moon near a node. Its
concept check asks why phase alone is insufficient. The export records
event\_prediction as false. Determining event occurrence, contacts and
local visibility requires further geometry; links lead to the existing
eclipse calculators without treating this diagram as an event
prediction.

The conceptual background is supported by NASA Science, Eclipses and the
Moon (https://science.nasa.gov/moon/eclipses/) and Moon Phases
(https://science.nasa.gov/moon/moon-phases/), consulted 4 November 2026.
Historical node and tithi references remain the
project\textquotesingle s Chapter 1 verses 43--44 and Chapter 2 verse
66. Modern explanatory references do not replace the adopted historical
model.

\textbf{4.7.4 Research Explorer}

The /research route exposes the saved longitude study for the Sun, Moon,
Mercury, Venus, Mars, Jupiter, Saturn and Rahu across eleven 100-year
windows from 1550 to 2649 CE. An 8 by 11 RMS matrix allows direct
selection; native body and window selectors provide the same choice.
Rahu is identified as the ascending lunar node. The two time treatments
share a vertical scale within each selected comparison.

\includegraphics[width=6.15in,height=3.42595in,alt={Figure 3 Time-aligned RMS longitude errors from the saved DE440 comparison}]{./media/current-report-3.png}

Figure 3 Time-aligned RMS longitude errors from the saved DE440
comparison

The naive treatment reads library JD directly as reference TT. The
time-aligned treatment subtracts 75.7873/360 day for the meridian and
adds the paper\textquotesingle s approximate Delta-T in days. It changes
the reference epoch, not the historical formulas. For the Moon in 1550
to 1649, the saved RMS changes from 3.990 to 1.686 degrees. A lower RMS
in one window is a diagnostic result, not calibration of the SS model.

The chart preserves the minimum and maximum sampled residual in each
five-year group and connects the group means. Full-window statistics use
every original ten-day sample. JSON and CSV downloads include source
paths and SHA-256 hashes, frame and time conventions, and sample counts.
The explorer does not need a live JPL request.

\textbf{5. Results And Findings}

\textbf{5.1 Worked-example case studies}

The inner-planet correction shows how a program can run with the wrong
interpretation. Treating Mercury's and Venus's revolution counts as
their own mean motion made the \emph{śīghra} anomaly zero and placed
Mercury 153 degrees from the Sun on one tested date. A bound on
elongation exposed the mistake. Another check fixed the sign of the
\emph{manda} anomaly: reversing apsis minus planet changes the
correction and can shift the Sun's computed place by several degrees.

A Kashi solar-eclipse example exposed the time-base error. With Ujjain
local time and Universal Time mixed, first-pass longitude parallax came
out at 3.83 \emph{ghaṭikās}, while Srivastava printed 1.06. A consistent
clock brought the computed value to 1.16. The remaining difference
matters because the commentary uses a more complete iteration and a
different disc-size treatment. These cases justify keeping intermediate
quantities and source rules visible, instead of judging the software by
a final answer alone.

Each case study identifies a symptom and its cause. For Mercury and
Venus, the symptom was an impossible elongation; the cause was confusing
a \emph{śīghrocca} revolution count with a mean planetary place. For
\emph{manda}, the symptom was a correction in the wrong direction; the
cause was reversing the textual anomaly convention. For the Kashi
eclipse, the symptom was a large parallax discrepancy; the cause was
mixing clocks. In each case a more precise test was added after the
reading changed. This makes the implementation a revisable account of
the text, rather than a collection of isolated outputs.

Another case involves the lunar apparent diameter. An initially
plausible scaling placed the largest disc at apogee; the verse-based
calculation scales the mean diameter by the ratio of true to mean daily
motion. The worked example's intermediate numbers exposed the reversal.
Chapter 12 provides a different kind of check: orbit values derived from
the shared-distance rule nearly reproduce the printed values. Chapter 14
tests can protect named seasons and years against off-by-one shifts.
Together, these cases show why verification needs more than one metric.
It uses domain bounds, arithmetic examples and cross-chapter consistency
where each is appropriate.

\textbf{5.2 DE440 evaluation}

A companion study compares Chapter 1--2 longitudes of the Sun, Moon,
five visible planets and lunar node with JPL DE440. It samples 40,176
epochs at ten-day intervals from 1550 to 2649 CE and considers both a
direct clock comparison and a Ujjain/ΔT-corrected variant. Independent
revalidation reconstructed all 642,816 stored residual values against
the local DE440 kernel to within about 0.00000055 degrees. This verifies
the arithmetic of the stored residuals under the stated frame and time
conventions. It does not validate every textual choice, later-chapter
event prediction or historical panchanga practice.

The comparison obtains tropical longitude by adding the
project\textquotesingle s ayanamsha to its sidereal longitude. A scalar
precession-corrected residual is a diagnostic on longitude, not a full
three-dimensional fixed-frame rotation. The text\textquotesingle s
trepidation model reaches a maximum near 2299 CE and then reverses,
affecting future extrapolation. The repository\textquotesingle s
revalidation note records concerns from an earlier manuscript snapshot;
it must be read in the context of its October 2026 review date and
revision. The current repository also supplies regenerate\_paper.py and
REPRODUCE.md. The new explorer does not rerun that pipeline or
independently reassess its fitted coefficients.

The sampled DE440 interval lies inside the local
kernel\textquotesingle s coverage. The study uses geometric geocentric
vectors, converts them to a mean ecliptic and equinox of date, and
compares wrapped longitude differences. The Moon\textquotesingle s node
uses a separate reference derived from its orbital angular momentum, so
its residual differs physically from a planet\textquotesingle s. Future
epochs test how the adopted constants and trepidation model extrapolate;
they are not observations already made. The current explorer groups the
saved decimal-year samples into the same windows as the saved tables,
with every exported group checked against those tables.

The earlier revalidation identified twelve epochs whose former
decimal-year labels assigned them to the preceding century. That finding
shows why residual arithmetic and subsequent grouping need separate
checks. The current explorer tests verify all 176 body-window-treatment
groups against the current saved tables, along with source hashes,
sample coverage and retained extrema. This establishes faithful
presentation of the supplied evidence, not an independent rerun of the
kernel or a validation of event timing. Angular elongation errors should
not be converted directly into claims about eclipse or tithi timing
without checking actual event crossings.

\textbf{5.3 Interface validation and reproducibility}

The repository verification record dated 8 November 2026 reports that
PYTHONPATH=src python -m pytest -q passed 1,677 tests after the corpus
repair. The learning-path validation record dated 4 November reports
1,653 passing tests, an earlier focused run of 49 tests across
test\_learning\_lab\_data.py, test\_research\_explorer\_data.py and
test\_api\_book\_ch02.py, and a successful production frontend build.
These results come from dated repository records; the tests were not
rerun while this report was edited.

\begin{longtable}[]{@{}
  >{\raggedright\arraybackslash}p{(\linewidth - 2\tabcolsep) * \real{0.5000}}
  >{\raggedright\arraybackslash}p{(\linewidth - 2\tabcolsep) * \real{0.5000}}@{}}
\toprule\noalign{}
\begin{minipage}[b]{\linewidth}\raggedright
Check
\end{minipage} & \begin{minipage}[b]{\linewidth}\raggedright
Evidence
\end{minipage} \\
\midrule\noalign{}
\endhead
\bottomrule\noalign{}
\endlastfoot
Source-data integrity & All 176 body-window-treatment groups match the
saved tables; each body has 40,176 samples per treatment. \\
Aggregation and provenance & Source hashes, complete sample counts, RMS,
signed means and sampled extrema are checked; generated files must match
fresh exports. \\
Browser workflows & Real API values, correction stages, whole-day
stepping, date-preserving links, rapid responses, failure/retry,
optional corpus, understanding feedback and exports. \\
Research interaction & 88 matrix cells; all 16 body/treatment
combinations; JSON and CSV metadata; keyboard selection and mobile
overflow checks. \\
Visual review & Light and dark desktop layouts and a 390-pixel mobile
layout; screenshot evidence retained with the project. \\
\end{longtable}

The frontend data exporters used for reproducibility are
docs/evaluation/export\_learning\_lab.py and
export\_research\_explorer.py. They summarize saved CSVs; they do not
compute fresh DE440 reference positions. The source-study pipeline
remains separately documented in docs/evaluation/REPRODUCE.md.

Browser checks in frontend cover all eight lessons and the research
explorer. They verify API values, lesson evidence exports, node
geometry, quizzes, date-preserving source links, stale responses, retry
and mobile layout. The calculation remains usable when corpus text is
missing. These checks establish the tested software behavior; historical
authenticity, eclipse timing accuracy and educational effectiveness
require separate evidence.

\textbf{6. Summary And Future Plan Of Work}

\textbf{6.1 Summary of work completed}

This project delivers a chapter-wise calculation library covering the
computable core of all fourteen chapters of the Sūrya Siddhānta,
together with its interfaces, source retrieval and evaluation. Its
strongest result is its record of how a textual rule becomes a
calculation: a passage fixes a rule, a function executes it, and a test
or worked example checks the consequence. This chain has corrected
plausible but wrong numbers. The interfaces make the work accessible,
while the DE440 comparison separately measures the adopted longitude
path against a modern reference. Known limits define what the software
currently supports. The completed learning paths and Research Explorer
let a user follow the work directly: choose an input, inspect correction
stages, follow the source reference and examine the comparison evidence.
They use the shared calculation library and saved study, with software
checks and explicit methodological limits.

\textbf{6.2 Limitations and next work}

The library covers the computable core of all fourteen chapters, not
every verse. Remaining gaps listed in the computational-edition
manuscript include full eclipse iteration, some disc-diameter and
visibility rules, parts of the equation of time and descriptive passages
that are not algorithms. The bundled Burgess corpus has complete
numbered-verse coverage, while remaining OCR noise and the optional
external commentary still require review. The interpretation of star
coordinates and procedures completed from commentary are documented
choices rather than uniquely settled historical facts. Modern
trigonometric operations also mean the software does not simulate every
rounding step of historical hand calculation.

Further work should archive the worked-example audit at a fixed code
revision, independently rerun and review the current DE440 regeneration
and fitting pipeline, and extend remaining procedures where a defensible
rule can be specified. The new teaching and research interfaces make
these questions visible, but do not resolve them. A learner study
remains future work. The eight completed learning paths cover the Sun,
Moon, five planets and Rahu, but the understanding checks alone do not
establish educational effectiveness.

The text is partly computational, partly descriptive. Completing every
verse as a function would require inventing a numerical rule for
passages that do not provide one. The project therefore makes a
narrower, more testable claim: it covers the computable core across all
chapters, with listed exceptions. Improvements should proceed through
explicit source decisions. For example, a full eclipse iteration can be
added with tests against the commentary's intermediate passes; a revised
star-coordinate reading can be compared against both current methods;
and the available DE440 pipeline can be independently rerun to
regenerate each table from a fixed kernel and code revision. The modular
structure allows these changes to be reviewed.

The remaining issues can be addressed in specific parts of the project.
The eclipse iteration belongs in the chapter-specific event functions
and can be checked against Srivastava's intermediate passes. The bundled
corpus needs continued OCR review and fixed-version provenance for
future revisions. The DE440 study needs independent rerunning and review
of the available regeneration and fitting pipeline, with a fixed code
and kernel revision, so another researcher can assess its figures and
fitted results. The star-coordinate alternatives need documented
comparisons before either is declared definitive. Addressing these
issues would make the project more useful for scholarly work by allowing
users to trace more conclusions from a stable source through executable
code to a repeatable result.

\textbf{References}

{[}1{]} Ebenezer Burgess, ``Translation of the Sūrya-Siddhānta, a
Text-Book of Hindu Astronomy; with Notes and an Appendix,'' Journal of
the American Oriental Society 6 (1860): 141--498.

{[}2{]} Mahavir Prasad Srivastava, Sūrya Siddhānta: Vigyan Bhashya, 2
vols., 2nd ed. (Allahabad: Dr. Ratnakumari Svadhyaya Sansthan,
1982--83).

{[}3{]} Kim Plofker, Mathematics in India (Princeton University Press,
2009).

{[}4{]} Ryan S. Park et al., ``The JPL Planetary and Lunar Ephemerides
DE440 and DE441,'' Astronomical Journal 161, no. 3 (2021): 105,
doi:10.3847/1538-3881/abd414.

{[}5{]} Pandu Santhoju et al., ``Mathematical Model as per the Algorithm
of Sūryasiddhānta for Computation of Location Specific True Position of
the Planets,'' Journal of Astrophysics and Astronomy 46 (2025): 9,
doi:10.1007/s12036-024-10018-8.

{[}6{]} Ashutosh Nigam, ``Evaluating the Accuracy of Surya Siddhanta
Planetary Algorithms against JPL DE440, 1550--2649 CE,'' unpublished
project manuscript (2026).

\end{document}
